\chapter{The Environment Movement}

\section{Environmental Racism and Eco-Justice}

\begin{KeyConcept}
\begin{itemize}
\item \textbf{Eco-justice} emerged in response to environmental inequities --- the \textbf{differential treatment of communities} regarding ecological concerns. Some communities bear a \textbf{disproportionate burden} of environmental problems.
\item \textbf{Environmental racism} (a concept of \textbf{Robert Bullard}): environmental policy/practice negatively affects groups based on class, race or ethnicity.
\item The \textbf{NIMBY syndrome} ('Not In My Backyard'): developed countries dump their waste elsewhere --- e.g. solid waste in America is disposed in \textbf{black/African-American settlements}, and Delhi's waste is dumped in unregulated settlements/slums. \textbf{Life chances are determined by geographical location, class and race}.
\end{itemize}
\end{KeyConcept}

\begin{ExampleBox}
\textbf{Eco-terrorism} (e.g. the \textbf{Earth Liberation Front}) is radical grassroots environmentalism that uses \textbf{violent and illegal methods} against businesses/organisations (mining etc.) believed to threaten the environment --- protecting the environment through non-legitimate means.
\end{ExampleBox}

\section{Indian vs Western Environmentalism}

\begin{center}
\begin{tabularx}{\textwidth}{YYY}
\rowcolor{AccentDark}
\thead{Aspect} & \thead{Western environmentalism} & \thead{Indian environmentalism} \\
Ideology & \textbf{Anthropocentric} --- human needs prioritised over nature (e.g. America withdrawing from Kyoto) & \textbf{Ecocentric} --- ecological needs prioritised \\
Religious basis & Consistent with \textbf{Judeo-Christian} belief in man's domination over earth; \textbf{human exceptionalism} & \textbf{Reverence for nature integral to Hindu philosophy} (Vayu dev, Agni dev, Ganga maa; Advaita Vedanta --- Adi Shankara's purusha--prakriti monism) \\
Origin & \textbf{Born of scientists} (IPCC, UNDP, Rio) & Born as a \textbf{response to the dominant notion of development/modernity} --- blossoming in the 1970s with Chipko and tribal resistance (e.g. the Dongria Kondh against bauxite mining) \\
Class interest & Represents the \textbf{upper and middle class} (`iPhone is eco-friendly', air purifiers) & The \textbf{environmentalism of the poor} \\
Motivation & \textbf{Leisure} --- keeping beautiful places unpolluted & \textbf{Survival} --- livelihood depends on the environment \\
Conservation ideology & \textbf{Protectionist conservationism} & \textbf{Utilitarian conservationism} (people depend on forest produce) \\
Measure & \textbf{GNP} (gross national product) & \textbf{Gross nature product} \\
\end{tabularx}
\end{center}

\begin{KeyConcept}
\textbf{Convergence --- Ulrich Beck's 'Risk Society'}: modern threats (COVID, global warming, bad air) affect \textbf{everyone}, not just a few sections --- \textbf{'wealth is hierarchic, but smog is democratic'} --- so the concern is no longer that of the poor alone (though the rich and poor manage the risk differently).
\end{KeyConcept}

\begin{QuickRevision}
\begin{itemize}
\item Western = anthropocentric, Judeo-Christian (man's domination), born of scientists (IPCC), upper/middle-class, leisure, protectionist, GNP.
\item Indian = ecocentric, Hindu reverence for nature (Advaita), response to development/modernity, environmentalism of the poor, survival, utilitarian, gross nature product.
\item Convergence (Beck): 'wealth is hierarchic, smog is democratic' --- modern risks affect everyone.
\end{itemize}
\end{QuickRevision}

\section{The Evolution of Modern Indian Environmentalism}

\begin{KeyConcept}
\begin{itemize}
\item The evolution of modern Indian environmentalism:
\item \textbf{1970s} --- the \textbf{'environmentalism of the poor'} (Ram Chandra Guha traces its origin to \textbf{Chipko}, followed by protest against commercial logging in the Himalayan foothills, involving both Gandhian and left-wing activists).
\item \textbf{1980s} --- institutionalisation: the \textbf{Ministry of Environment and Forests (MoEF\&CC)} was set up (upgraded from a department), followed by state-level departments.
\item \textbf{1990s} --- every university made an \textbf{environment course mandatory}; the fields of ecological, economic and environmental history emerged.
\item \textbf{1995} --- \textbf{anti-environmental backlash}: the Greens, who had made the state accountable (asking where land for new projects would come from), were regarded as the \textbf{only obstacle to India's march to greatness}.
\end{itemize}
\end{KeyConcept}

\begin{KeyConcept}
\begin{itemize}
\item \textbf{India's environmental-management policy followed the 'frontier economics model'} --- treating the environment as an \textbf{infinite supply of physical resources} for human society (nature--society relation dominated by humans). What is needed instead is an \textbf{'eco-development model'} in which development is \textbf{synergetic with ecosystem processes} --- moving from the \textbf{'polluter pays' principle} to the \textbf{'pollution prevention pays' concept}.
\item The emergence of these movements paved the way for a new development paradigm --- \textbf{sustainable development} (the \textbf{new ecological paradigm}).
\end{itemize}
\end{KeyConcept}

\section{Guha and Gadgil: Eight Environmentalists and Three Individuals}

\begin{KeyConcept}
\begin{itemize}
\item \textbf{Ram Chandra Guha and Madhav Gadgil} (in \textbf{'Ecology and Equity'}) described \textbf{eight types of environmentalists} in India:
\item 1. \textbf{Crusading Gandhians} --- cite Hindu scriptures; advocate self-realization, ecological prudence, frugality and a return to the pre-capitalist village; against mines, power projects and railways; use non-violent methods (Chipko, Narmada Bachao).
\item 2. \textbf{Ecological Marxists} --- see unequal access to resources as the cause of degradation; favour local-level decision-making (people's science movements); may use Gandhian tactics (the \textbf{Silent Valley} movement in Kerala --- which got a hydroelectric project abandoned).
\item 3. \textbf{Appropriate technologists} --- seek a synthesis of agriculture and industry; advocate small-scale, intermediate, eco-friendly technology (the appropriate-technology wing of Chipko, organising soil/water conservation and afforestation).
\item 4. \textbf{Wilderness enthusiasts} --- e.g. BBC's Planet Earth (David Attenborough); Sunderlal Bahuguna; Greta Thunberg.
\item 5. \textbf{Guardians of the sacred} --- protecting sacred groves (North-East and tribal groups).
\item 6. \textbf{Community managers} --- village forest committees managing 'safety forests' (sacred groves) and 'supply forests' (regulated harvest).
\item 7. \textbf{Technocrat planners} --- who carry out the \textbf{Environmental Impact Assessment (EIA)}.
\item 8. \textbf{Judicial activists} --- the National Green Tribunal, and public interest litigation (PIL).
\end{itemize}
\end{KeyConcept}

\begin{KeyConcept}
Gadgil and Guha also observed that Indian development policies create \textbf{three types of individuals}: the \textbf{omnivores} (the elite beneficiaries of development projects), the \textbf{ecosystem people} (those dependent on the natural environment --- tribals), and the \textbf{ecological refugees} (displaced by mining/development). Independent India has become a \textbf{field of conflict between these groups}, triggered by the abuse of natural resources to benefit the omnivore elite.
\end{KeyConcept}

\begin{KeyConcept}
\begin{itemize}
\item \textbf{Anil Agarwal} described \textbf{four rays of hope} for India's environment: the \textbf{NGO movement, the judiciary, people's initiatives} (e.g. Ralegan Siddhi --- Anna Hazare's village) and the \textbf{media}. He called for a \textbf{'third war of independence'} --- to be fought by post-independence Indians \textbf{with themselves} to protect their environment. \textbf{Dr. M. P. Parameswaran} spoke of building a \textbf{'fourth world'} --- combining Marxism, Gandhism, environmentalism, eco-feminism and human rights.
\item \textbf{Akhileshwar Pathak} argues that the Indian state largely promoted modernization and capitalism --- so the poor, in order to survive, had to collaborate with their oppressors.
\end{itemize}
\end{KeyConcept}

\section{The Chipko Movement}

\begin{KeyConcept}
\begin{itemize}
\item \textbf{Chipko} (1970s) was an organised resistance to the destruction of forests: ash trees were allotted to private companies (for cement), against which the \textbf{Dasholi Gram Swarajya Mandal} (a local cooperative) mobilised. Women \textbf{lay down in front of the trucks} (non-violent, Gandhian), and \textbf{Chandi Prasad Bhatt} suggested \textbf{embracing (hugging) the trees} to prevent them being cut.
\item Its \textbf{six demands}: (1) recognition of the \textbf{traditional rights} of forest people over forest produce; (2) making arid forests green through \textbf{people's participation}; (3) forming \textbf{village committees} to manage forests; (4) developing \textbf{forest-based home industries} (raw material, money and technique); (5) giving priority to \textbf{local varieties} in afforestation; and (6) asking men to \textbf{give up alcohol}.
\end{itemize}
\end{KeyConcept}

\section{Theories: Neo-Marxist Environmental Sociology}

\begin{KeyConcept}
\begin{itemize}
\item The \textbf{neo-Marxist theory} holds that development comes at the cost of natural resources --- so capitalism is not sustainable; development and nature are opposed.
\item \textbf{Treadmill of production}: the modern economic system requires \textbf{constant growth} (the iPhone 11, 12, 13...), causing increased resource exploitation; in the 20th century, \textbf{consumerism} (the middle/leisure class displaying status through goods) keeps the treadmill running.
\item \textbf{Ecological Marxism} (John Bellamy Foster): 'labour is the father of material wealth, \textbf{earth is its mother}' --- the force that subjugates the proletariat is the same force that destroys the earth. Marx condemned capitalism not only for exploiting men and women but for the \textbf{metabolic rift} (his term) --- the imbalance between human needs and nature's carrying capacity.
\item Capitalism \textbf{privatises the commons} (forest, air, water; even the 'spectrum'/air is auctioned), and rests on an \textbf{inherent contradiction}: more machines to cut labour costs means more exploitation of nature --- 'workers of the world, you have nothing to lose but the environment'.
\item \textbf{World-systems theory}: the world is divided into core and periphery; the rich nations are the main consumers but the \textbf{environmental damage falls heaviest on the poor nations} (the concept of \textbf{historical responsibility} in the Kyoto Protocol).
\end{itemize}
\end{KeyConcept}

\section{Ecological Modernization and Its Critique}

\begin{KeyConcept}
\begin{itemize}
\item The \textbf{neoliberal theory --- ecological modernization theory} --- holds that \textbf{modernization is compatible with ecological sustainability} (there is no conflict): the way out of the ecological crisis is \textbf{further industrialisation} (green growth/green capitalism --- catalytic converters, BS6 emission standards, e-vehicles, solar tracking, vertical gardens).
\item It supports \textbf{green consumerism and decarbonization}, and informs policy (green contracts, carbon tax, green GDP, the Kyoto Protocol). Its \textbf{Environmental Kuznets Curve} says development first degrades, then protects, the environment (an inverted-U: after the turning point, environment improves).
\item \textbf{Critiques}: \textbf{Jevons paradox} (fuel-efficiency gains increase, not decrease, fuel use --- cheaper driving means more people drive); the \textbf{Netherlands fallacy} (wrongly assuming a rich nation's environmental impact is contained within its national borders).
\end{itemize}
\end{KeyConcept}

\section{Ulrich Beck's Risk Society}

\begin{KeyConcept}
\begin{itemize}
\item \textbf{Ulrich Beck ('Risk Society')}: traditional societies faced \textbf{external risks} (drought, earthquake, famine --- unrelated to human action, affecting few sections); modern societies face \textbf{manufactured risks} --- the outcomes of our \textbf{own intervention into nature} (nuclear holocaust, biological weapons, laboratory errors like COVID).
\item \textbf{Martin Rees} ('Our Final Century') argued we must focus on the \textbf{unintended consequences of scientific advancement}. As environmental risks rise, the old class conflicts over the 'wealth cake' lose significance --- because the cake itself is poisoned.
\item We are thus moving into a \textbf{world risk society}, in which \textbf{risk consciousness and risk avoidance} become central --- because pollution does not respect national boundaries (hence the idea of the global commons).
\item In the pandemic, the virus was 'the most democratic virus' --- a manufactured risk affecting everyone.
\end{itemize}
\end{KeyConcept}

\begin{QuickRevision}
\begin{itemize}
\item Environmental racism (Bullard); NIMBY; eco-terrorism (Earth Liberation Front).
\item Indian (ecocentric, survival, poor, utilitarian, gross-nature-product) vs Western (anthropocentric, leisure, upper-middle class, protectionist, GNP) environmentalism; convergence via Beck ('smog is democratic').
\item Guha \& Gadgil: 8 environmentalists (Crusading Gandhians, Ecological Marxists, Appropriate technologists, Wilderness enthusiasts, Guardians of the sacred, Community managers, Technocrat planners, Judicial activists); 3 individuals (omnivores, ecosystem people, ecological refugees).
\item Chipko: Dasholi Gram Swarajya Mandal, Chandi Prasad Bhatt (tree-hugging); six demands.
\item Neo-Marxist theory: treadmill of production; ecological Marxism (Foster; 'earth is its mother'; metabolic rift); privatisation of commons; world-systems (core-periphery, historical responsibility).
\item Ecological modernization: green capitalism, Environmental Kuznets Curve; critiques --- Jevons paradox, Netherlands fallacy.
\item Beck's Risk Society: external vs manufactured risks; Martin Rees; world risk society.
\item PYQs: Chipko as eco-feminism; Western patriarchy (Vandana Shiva); 'has green-cover reduction caused global warming' (Neo-Marxist yes / eco-modernization no); 'anthropogenic influence on climate' (2020); 'humanity at the mercy of nature/technology' (2020) --- conclude with the new ecological paradigm (polluter pays $\rightarrow$ pollution prevention; material-centric $\rightarrow$ human-centric/Amartya Sen's capability approach).
\end{itemize}
\end{QuickRevision}

